\section{激活函数与门控线性单元}

\subsection{从 ReLU 到 GELU}

\begin{figure}[H]
\centering
\includegraphics[width=0.95\textwidth]{slides-images/slide-016.jpg}
\caption{激活函数一览：ReLU、GELU、Swish 以及各种 GLU 变体\protect\footnotemark}
\end{figure}
\footnotetext{Slides 第17页。}

激活函数的演变：

\begin{itemize}
    \item \textbf{ReLU}：$\text{ReLU}(x) = \max(0, x)$，原始 Transformer 使用
    \item \textbf{GELU}：$\text{GELU}(x) = x \cdot \Phi(x)$，其中 $\Phi$ 是标准正态分布的 CDF。GPT 系列使用
    \item \textbf{Swish/SiLU}：$\text{Swish}(x) = x \cdot \sigma(x)$，其中 $\sigma$ 是 sigmoid 函数
\end{itemize}

GELU 和 Swish 的形状类似 ReLU，但在零点附近更平滑（可微），且在负值区域有一个小的``凹陷''。

\subsection{门控线性单元（GLU）}

\begin{figure}[H]
\centering
\includegraphics[width=0.95\textwidth]{slides-images/slide-017.jpg}
\caption{标准 FFN vs 门控线性单元（GLU）：GLU 在隐藏层添加了一个门控分支\protect\footnotemark}
\end{figure}
\footnotetext{Slides 第18页。}

门控线性单元是现代 Transformer 中最重要的架构创新之一。标准 FFN：
$$\text{FFN}(x) = \text{activation}(xW_1)W_2$$

GLU 变体：
$$\text{GLU-FFN}(x) = (\text{activation}(xW_1) \odot xV) W_2$$

\begin{itemize}
    \item $W_1$：第一个线性变换（投影到隐藏空间）
    \item $V$：门控矩阵（额外参数）
    \item $\odot$：逐元素乘法（门控操作）
    \item $W_2$：第二个线性变换（投影回原始空间）
\end{itemize}

\begin{importantbox}{三种 GLU 变体}
\begin{itemize}
    \item \textbf{ReGLU}：$(\text{ReLU}(xW_1) \odot xV) W_2$
    \item \textbf{GEGLU}：$(\text{GELU}(xW_1) \odot xV) W_2$，Google T5V1.1、Gemma 系列使用
    \item \textbf{SwiGLU}：$(\text{Swish}(xW_1) \odot xV) W_2$，LLaMA、PaLM、OLMo 等使用，\textbf{是当前最流行的选择}
\end{itemize}
\end{importantbox}

\begin{figure}[H]
\centering
\includegraphics[width=0.95\textwidth]{slides-images/slide-019.jpg}
\caption{GLU 变体的实验对比：GLU 变体（加粗）在各项指标上一致优于非门控变体\protect\footnotemark}
\end{figure}
\footnotetext{Slides 第20页。来自 Shazeer 2020 的原始 GLU 论文。}

\begin{figure}[H]
\centering
\includegraphics[width=0.95\textwidth]{slides-images/slide-020.jpg}
\caption{Narang et al. 2020 的消融实验：GLU 变体一致取得更低的损失\protect\footnotemark}
\end{figure}
\footnotetext{Slides 第21页。}

\begin{knowledgebox}{门控机制的普遍有效性}
门控（Gating）是深度学习中反复出现的有效模式。从 LSTM 的门控机制到 Highway Networks，再到 Transformer 中的 GLU，门控操作让网络能够动态选择性地传递信息。GLU 本质上是让网络学习``哪些隐藏特征应该被保留，哪些应该被抑制''。
\end{knowledgebox}

\begin{warningbox}{GLU 不是必须的}
虽然 GLU 变体普遍更好，但不使用 GLU 也能训练出优秀的模型。例如：
\begin{itemize}
    \item GPT-3 使用标准 GELU（非门控），仍是经典模型
    \item Nemotron 340B 使用 Squared ReLU
    \item Falcon 2 11B 使用标准 ReLU
\end{itemize}
GLU 提供了一致的小幅增益，但绝非成功的必要条件。
\end{warningbox}

\subsection{本章小结}

\begin{itemize}
    \item SwiGLU 是当前最流行的 FFN 激活方式，在保持参数量不变的前提下一致优于非门控变体
    \item 门控线性单元引入了额外的参数矩阵 $V$，为保持总参数量不变，通常将隐藏维度缩小为原来的 $2/3$
    \item 激活函数的具体选择（ReLU vs GELU vs Swish）对性能的影响远小于是否使用门控机制
\end{itemize}

%% ========== Part 4: 串并行层与位置编码 ==========

